\clearpage
\fancyhead[L]{\small GAUSS AN BODE.}
\hypertarget{native-298-lower}{}
\begin{center}
GAUSS AN BODE.
\end{center}

\begin{center}
\begin{tabular}{c}
\hline
Astronomisches Jahrbuch für 1811. S. 135--139. Berlin 1808.\\
\hline
\end{tabular}
\end{center}

\begin{flushright}
Göttingen, 24. Januar 1808.
\end{flushright}

{\rightskip=0pt plus 1em\par\indent Zunächst werden die mit meinem Ortswechsel; verbundenen Unruhen und danach meine durch äußerst ungünstiges Wetter vereitelte Hoffnung, Ihre Kometenbeobachtungen mit zahlreichen hiesigen Beobachtungen zu erwidern, mein langes Schweigen entschuldigen. Hier erhalten Ew. --- unterdessen die wenigen Beobachtungen des Kometen, die ich dem hier fast ununterbrochen bedeckten Himmel habe abgewinnen können.\par}

\begin{center}\fontsize{9}{11}\selectfont\setlength{\tabcolsep}{2pt}
\begin{tabular}{lrrrl}
& M. Z. in Göttingen. & Rektasz. & Dekl. & \\
1807. Dez. 17. & $11^{\text{u}}\,24^{\text{m}}\,44^{\text{s}}$ & $305^\circ\;54'\;0''$ & $45^\circ\;4'\;30''$ & \\
29. & $7\;59\;21$ & $318\;57\;41$ & $46\;41\;48$ & \\
& $11\;36\;33$ & $319\;6\;35$ & $46\;43\;11$ & \\
1808. Jan. 4. & $7\;22\;12$ & $325\;11\;19$ & $47\;12\;46$ & \smash{\raisebox{16.5pt}{$\left.\vphantom{\begin{array}{c}0\\0\\0\\0\end{array}}\right\}$ von mir.}}\\
& $9\;13\;57$ & $325\;15\;53$ & & \\
& $17\;30\;26$ & $325\;36\;38$ & & von Prof. \textsc{Harding}.\\
\end{tabular}
\end{center}

\par\indent Die parabolische Bahn des Kometen hatte ich bereits Anfang November aus einigen Bremer und Lilienthaler Beobachtungen berechnet. Da diese bisher noch nicht bekannt geworden sind, teile ich sie Ihnen hier mit.

\SourceLeaf{299}\noindent Durchgang durch das Perihelium 18. Sept. $19^{\text{u}}\,16^{\text{m}}\,0^{\text{s}}$ Merid. von Paris.


\begin{center}\fontsize{10}{12}\selectfont\setlength{\tabcolsep}{3pt}\renewcommand{\arraystretch}{1.15}\begin{tabular}{lr}
Länge des Perihels .\;.\;.\;.\;.\;.\;.\;.\;.\;.\;.\;.\;.\; & $271^\circ\;0'\;13''3.$ \\
Logarithmus des kleinsten Abstands , .\;.\;.\;.\;.\;.\;.\;.\;.\;.\; & $9.8114927$ \\
Aufsteigender Knoten .\;.\;.\;.\;.\;.\;.\;.\;.\;.\;.\;.\;.\; & $266^\circ\;38'\;31''2.$ \\
Neigung der Bahn .\;.\;.\;.\;.\;.\;.\;.\;.\;.\;.\;.\;.\;.\; & $63\;12\;35.6.$ \\
\end{tabular}\end{center}

\noindent Bewegung rechtläufig.


\par\indent Ende November und Anfang Dezember wichen diese Elemente um etwa $1'$ von den Beobachtungen ab, danach nahm der Fehler aber wieder ab. Sollte es mir gelingen, noch einige Beobachtungen des Kometen zu erhalten, werde ich die oben angegebene Bahn danach noch einmal verbessern. Die Korrekturen werden aber gewiss sehr klein sein. Über die Abweichung der Bahn von der Parabel, die nicht beträchtlich sein kann, wird sich anhand der Beobachtungen schwerlich etwas entscheiden lassen.

\par\indent Meine neuen (X.) Elemente der \emph{Pallas} sind folgende:


\begin{center}\fontsize{10}{12}\selectfont\setlength{\tabcolsep}{3pt}\renewcommand{\arraystretch}{1.15}\begin{tabular}{lr}
Epoche 1808 Meridian von Göttingen .\;.\;.\;.\;.\;.\;.\; & $252^\circ\;32'\;28''5$ \\
Tägliche tropische Bewegung .\;.\;.\;.\;.\;.\;.\;.\;.\;.\;.\;.\; & $770''2143$ \\
Logarithmus der großen Halbachse .\;.\;.\;.\;.\;.\;.\;.\;.\; & $0.4423149$ \\
Exzentrizität .\;.\;.\;.\;.\;.\;.\;.\;.\;.\;.\;.\;.\;.\;.\;.\;.\; & $0.2450198$ \\
Perihel 1803 .\;.\;.\;.\;.\;.\;.\;.\;.\;.\;.\;.\;.\;.\;.\; & $121^\circ\;3'\;11''4$ \\
$\AscendingNode\;1803$ .\;.\;.\;.\;.\;.\;.\;.\;.\;.\;.\;.\;.\;.\;.\;.\;.\;.\; & $172\;28\;56.9$ \\
Neigung der Bahn .\;.\;.\;.\;.\;.\;.\;.\;.\;.\;.\;.\;.\;.\; & $34\;37\;41.0$ \\
\end{tabular}\end{center}


\par\indent Die Elemente der \emph{Juno} habe ich vor Kurzem anhand der Beobachtungen verbessert, die Herr \textsc{Bessel} im vergangenen Jahr angestellt hat. Hier sind die neuen (VII.) Elemente:


\begin{center}\fontsize{10}{12}\selectfont\setlength{\tabcolsep}{3pt}\renewcommand{\arraystretch}{1.15}\begin{tabular}{lr}
Epoche, Meridian von Göttingen 1805 .\;.\;.\;.\;.\;.\;.\; & $42^\circ\;37'\;3''7$ \\
1806 .\;.\;.\;.\;.\;.\;.\; & $125\;7\;57.8$ \\
1807 .\;.\;.\;.\;.\;.\;.\; & $207\;38\;51.8$ \\
1808 .\;.\;.\;.\;.\;.\;.\; & $290\;23\;19.8$ \\
Mittlere tägliche tropische Bewegung .\;.\;.\;.\;.\;.\;.\;.\;.\; & $813''8468$ \\
Tropische Umlaufszeit .\;.\;.\;.\;.\;.\;.\;.\;.\;.\;.\; & $1592\;\text{Tage}\;10\tfrac12\;\text{St.}$ \\
Sonnennähe 1805 .\;.\;.\;.\;.\;.\;.\;.\;.\;.\;.\;.\;.\;.\;.\; & $53^\circ\;19'\;0''2$ \\
Aufsteigender Knoten 1805 .\;.\;.\;.\;.\;.\;.\;.\;.\;.\;.\; & $171\;4\;28.2$ \\
Exzentrizität .\;.\;.\;.\;.\;.\;.\;.\;.\;.\;.\;.\;.\;.\;.\;.\;.\; & $0.2554996$ \\
Logarithmus der großen Halbachse .\;.\;.\;.\;.\;.\;.\;.\;.\; & $0.4263781$ \\
Neigung der Bahn .\;.\;.\;.\;.\;.\;.\;.\;.\;.\;.\;.\;.\;.\; & $13^\circ\;4'\;26.2$ \\
\end{tabular}\end{center}


\par\indent Die daraus berechnete Ephemeride für 1808 lautet: [Sie gibt für Mitternacht in Göttingen vom 16. April bis zum 28. Dezember in Abständen von vier Tagen Rektaszension und Deklination auf Minuten genau an.]

\par\indent Mein Werk über die Bestimmung der Bahnen der Himmelskörper ist jetzt im Druck. Ich hoffe, dass es zu Ostern erscheinen kann.

